% Options for packages loaded elsewhere
\PassOptionsToPackage{unicode}{hyperref}
\PassOptionsToPackage{hyphens}{url}
\documentclass[
  polish,
  11pt,
  a4paper,
]{article}
\usepackage{xcolor}
\usepackage[margin=2cm]{geometry}
\usepackage{amsmath,amssymb}
\setcounter{secnumdepth}{-\maxdimen} % remove section numbering
\usepackage{iftex}
\ifPDFTeX
  \usepackage[T1]{fontenc}
  \usepackage[utf8]{inputenc}
  \usepackage{textcomp} % provide euro and other symbols
\else % if luatex or xetex
  \usepackage{unicode-math} % this also loads fontspec
  \defaultfontfeatures{Scale=MatchLowercase}
  \defaultfontfeatures[\rmfamily]{Ligatures=TeX,Scale=1}
\fi
\usepackage{lmodern}
\ifPDFTeX\else
  % xetex/luatex font selection
\fi
% Use upquote if available, for straight quotes in verbatim environments
\IfFileExists{upquote.sty}{\usepackage{upquote}}{}
\IfFileExists{microtype.sty}{% use microtype if available
  \usepackage[]{microtype}
  \UseMicrotypeSet[protrusion]{basicmath} % disable protrusion for tt fonts
}{}
\makeatletter
\@ifundefined{KOMAClassName}{% if non-KOMA class
  \IfFileExists{parskip.sty}{%
    \usepackage{parskip}
  }{% else
    \setlength{\parindent}{0pt}
    \setlength{\parskip}{6pt plus 2pt minus 1pt}}
}{% if KOMA class
  \KOMAoptions{parskip=half}}
\makeatother
\usepackage{longtable,booktabs,array}
\newcounter{none} % for unnumbered tables
\usepackage{calc} % for calculating minipage widths
% Correct order of tables after \paragraph or \subparagraph
\usepackage{etoolbox}
\makeatletter
\patchcmd\longtable{\par}{\if@noskipsec\mbox{}\fi\par}{}{}
\makeatother
% Allow footnotes in longtable head/foot
\IfFileExists{footnotehyper.sty}{\usepackage{footnotehyper}}{\usepackage{footnote}}
\makesavenoteenv{longtable}
\usepackage{graphicx}
\makeatletter
\newsavebox\pandoc@box
\newcommand*\pandocbounded[1]{% scales image to fit in text height/width
  \sbox\pandoc@box{#1}%
  \Gscale@div\@tempa{\textheight}{\dimexpr\ht\pandoc@box+\dp\pandoc@box\relax}%
  \Gscale@div\@tempb{\linewidth}{\wd\pandoc@box}%
  \ifdim\@tempb\p@<\@tempa\p@\let\@tempa\@tempb\fi% select the smaller of both
  \ifdim\@tempa\p@<\p@\scalebox{\@tempa}{\usebox\pandoc@box}%
  \else\usebox{\pandoc@box}%
  \fi%
}
% Set default figure placement to htbp
\def\fps@figure{htbp}
\makeatother
\ifLuaTeX
\usepackage[bidi=basic,shorthands=off]{babel}
\else
\usepackage[bidi=default,shorthands=off]{babel}
\fi
\ifLuaTeX
  \usepackage{selnolig} % disable illegal ligatures
\fi
\setlength{\emergencystretch}{3em} % prevent overfull lines
\providecommand{\tightlist}{%
  \setlength{\itemsep}{0pt}\setlength{\parskip}{0pt}}
\usepackage{fontspec}
\defaultfontfeatures{Ligatures=TeX}
\usepackage{amsmath,amssymb}
\usepackage{booktabs}
\usepackage{array}
\usepackage{geometry}
\usepackage{enumitem}
\usepackage{xcolor}
\usepackage{tcolorbox}
\usepackage{fancyhdr}
\usepackage{titlesec}
\usepackage{multicol}

\geometry{margin=1.8cm, top=2.5cm, bottom=2cm}
\pagestyle{fancy}
\fancyhf{}
\rhead{\small Projektowanie badań — 2026/27}
\lhead{\small Zarządzanie informacją | ISI UJ}
\cfoot{\thepage}

\titleformat{\section}{\large\bfseries\color{blue!70!black}}{\thesection.}{0.5em}{}
\titleformat{\subsection}{\normalsize\bfseries}{\thesubsection.}{0.5em}{}
\titleformat{\subsubsection}{\small\bfseries\color{gray!70!black}}{\thesubsubsection.}{0.5em}{}
\titlespacing*{\section}{0pt}{1.5ex}{1ex}
\titlespacing*{\subsection}{0pt}{1ex}{0.5ex}

\setlist[itemize]{nosep, leftmargin=1.5em}
\setlist[enumerate]{nosep, leftmargin=1.5em}

\tcbset{boxrule=0.5pt, left=4pt, right=4pt, top=4pt, bottom=4pt, fonttitle=\bfseries\small}

\newtcolorbox{definicja}[1][]{colback=blue!5, colframe=blue!50!black, title={#1}}
\newtcolorbox{uwagabox}[1][]{colback=orange!10, colframe=orange!70!black, title={#1}}
\newtcolorbox{wzorbox}[1][]{colback=gray!5, colframe=gray!50!black, title={#1}}
\newtcolorbox{przykladbox}[1][]{colback=purple!5, colframe=purple!50!black, title={#1}}
\usepackage{bookmark}
\IfFileExists{xurl.sty}{\usepackage{xurl}}{} % add URL line breaks if available
\urlstyle{same}
% fallback for those not using the hyperref driver hyperxmp:
\makeatletter
\@ifundefined{xmpquote}{\newcommand{\xmpquote}[1]{#1}}{}
\makeatother
\hypersetup{
  pdftitle={Handout W05 --- Korelacja{,} regresja i język raportu},
  pdfauthor={Marek Deja},
  pdflang={pl-PL},
  hidelinks,
  pdfcreator={LaTeX via pandoc}}

\title{Handout W05 --- Korelacja, regresja i język raportu}
\usepackage{etoolbox}
\makeatletter
\providecommand{\subtitle}[1]{% add subtitle to \maketitle
  \apptocmd{\@title}{\par {\large #1 \par}}{}{}
}
\makeatother
\subtitle{Mapa pojęć i sprawdziany rozumienia}
\author{Marek Deja}
\date{2026/27}

\begin{document}
\maketitle

\section{Co chcemy wiedzieć?}\label{co-chcemy-wiedzieux107}

W S02 x oznacza deklarowaną użyteczność katalogu, a y obserwowany czas.
Indeks nie jest testem całej kompetencji. Ujemny związek oznacza
współwystępowanie wyższej oceny z krótszym czasem; krótki czas
porzuconego zadania nie jest sukcesem.

Jednostką jest osoba z kompletną parą. Pytanie o korelację, średnią
przewidywaną i zgodność pomiarów wymaga różnych odpowiedzi. Przed
obliczeniem ustal konstrukt, skalę, populację i regułę braków.

{\def\LTcaptype{none} % do not increment counter
\begin{longtable}[]{@{}
  >{\raggedright\arraybackslash}p{(\linewidth - 4\tabcolsep) * \real{0.3333}}
  >{\raggedright\arraybackslash}p{(\linewidth - 4\tabcolsep) * \real{0.3333}}
  >{\raggedright\arraybackslash}p{(\linewidth - 4\tabcolsep) * \real{0.3333}}@{}}
\toprule\noalign{}
\begin{minipage}[b]{\linewidth}\raggedright
Pytanie
\end{minipage} & \begin{minipage}[b]{\linewidth}\raggedright
Estymanda
\end{minipage} & \begin{minipage}[b]{\linewidth}\raggedright
Wynik próbny
\end{minipage} \\
\midrule\noalign{}
\endhead
\bottomrule\noalign{}
\endlastfoot
Jak silna jest liniowa współzmienność? & \(\rho_P\) & r Pearsona \\
Jak zgodna jest kolejność? & \(\rho_S\) & \(r_S\) Spearmana \\
O ile różni się średni y na jednostkę x? & \(\beta_1\) & nachylenie
\(b_1\) \\
Czy samoocena trafnie podaje poziom? & np. średni błąd w porównywalnej
skali & różnice, nie samo r \\
\end{longtable}
}

\section{Przeczytaj wzór zdaniem}\label{przeczytaj-wzuxf3r-zdaniem}

\[
r=\frac{\sum(x_i-\bar x)(y_i-\bar y)}
{\sqrt{\sum(x_i-\bar x)^2\sum(y_i-\bar y)^2}}.
\]

Licznik łączy odchylenia obu zmiennych od ich średnich; mianownik
standaryzuje wynik. r jest bez jednostki, ma zakres {[}−1; 1{]}, opisuje
liniowość. Zero nie wyklucza zależności w kształcie U. Spearman jest
korelacją rang, przy remisach średnich rang. Wybór metody nie sprowadza
się do wyniku testu normalności.

\[
t=r\sqrt{\frac{n-2}{1-r^2}},\quad df=n-2.
\]

To klasyczny test zerowej korelacji Pearsona; model niezależnych par
dwuwymiarowo normalnych uzasadnia rozkład t. Dla wyniku Spearmana
odczytujemy właściwy wynik programu, nie dopisujemy fikcyjnych df do
statystyki S.

\section{Dwa losowania, dwa pytania}\label{dwa-losowania-dwa-pytania}

W permutacji przestawiamy y względem x. Zachowujemy rozkłady, zmieniamy
parowanie pod hipotezą niezależności i wymienialności. W Monte Carlo
dwustronne p ma postać \((k+1)/(B+1)\), gdzie k liczy \(|r^*|\geq|r|\).
W pełnej enumeracji dzielimy liczbę skrajnych układów przez wszystkie
układy bez poprawki +1.

W bootstrapie losujemy całe pary ze zwracaniem. Zachowujemy relację
wewnątrz osoby, aby przybliżyć niepewność współczynnika. Oddzielne
losowanie kolumn zniszczyłoby tę relację. Raportujemy metodę CI i liczbę
prawidłowych replik.

Swobodne tasowanie nie uwzględnia automatycznie powtarzanych pomiarów
ani klastrów. Zerowa korelacja i niezależność nie są równoważne dla
wszystkich rozkładów. Zgodne p nie usuwają problemu pomiaru lub doboru
próby.

\section{Czytamy model}\label{czytamy-model}

\[
E(Y\mid X=x)=\beta_0+\beta_1x,\qquad
\widehat y=b_0+b_1x,\qquad e_i=y_i-\widehat y_i.
\]

\(b_1\) ma jednostkę y na jednostkę x, tutaj minuty/punkt. \(b_0\)
dotyczy x=0, które może leżeć poza skalą. Reszta jest różnicą wyniku
zaobserwowanego i przewidywanego, nie błędem wpisania danych.

\[
R^2=1-\frac{\sum e_i^2}{\sum(y_i-\bar y)^2}.
\]

Dla OLS z wyrazem wolnym jest to część sumy kwadratów zmienności y
odtworzona przez model w próbie. W regresji prostej na tych samych
parach \(R^2=r_P^2\), nie \(r_S^2\). Nie oznacza procentu poprawnie
sklasyfikowanych osób ani siły dowodu przyczynowego.

\section{Mały przykład
wykładowy}\label{maux142y-przykux142ad-wykux142adowy}

Sześć par: x=(1,2,2,3,4,5), y=(15,12,14,10,9,6). To celowo regularne
dane dydaktyczne, inne niż pełny zbiór C08.

Pearson: r≈−0,974; t(4)≈−8,668; p≈0,000975. Wszystkie 720 permutacji
indeksów dają 4 wyniki co najmniej tak skrajne: p\_perm≈0,00556. Różne
procedury korzystają z różnych rozkładów odniesienia.

Model: \(\widehat y=17{,}277-2{,}215x\). Nachylenie około −2,22
min/punkt, 95\% CI {[}−2,93; −1,51{]}. To różnica przewidywanej średniej
między poziomami x, nie efekt podniesienia czyjejś odpowiedzi
ankietowej.

Dla x=3: prognoza 10,63 min, CI średniej {[}9,67; 11,59{]}, przedział
nowej osoby {[}8,11; 13,16{]}. Szerszy przedział zawiera dodatkową
zmienność indywidualną. Nie ekstrapolujemy na x=8 na skali 1--5.

\section{Co sprawdzamy na wykresie?}\label{co-sprawdzamy-na-wykresie}

Krzywiznę, rozrzut zależny od poziomu, skupienia, nietypowe punkty i
zakres. Dźwignia dotyczy nietypowego x, reszta odchylenia y od prostej,
wpływ zmiany dopasowania. Punktu nie usuwa się tylko dlatego, że zmienia
p.

Klasyczne przedziały regresji wymagają założeń o średniej, niezależności
i wariancji błędów; dokładne małopróbkowe wnioskowanie t również ich
normalności. Test normalności nie sprawdza reprezentatywności ani
trafności narzędzia.

\section{Trzy zdania do poprawienia}\label{trzy-zdania-do-poprawienia}

„r=−0,30 to spadek kompetencji o 30\%''. Korelacja nie jest procentem
zmiany, a użyteczność katalogu nie jest całym konstruktem kompetencji.

„r=1 dowodzi trafnej samooceny''. Można idealnie zachować kolejność i
zawyżać każdy wynik o 20 punktów procentowych. Korelacja nie mierzy
zgodności poziomów.

„CI obejmuje zero, więc nie ma żadnego związku''. Przedział pokazuje
niepewność konkretnego parametru; dodatkowo liniowa miara może pomijać
nieliniowy wzorzec.

\section{Akapit wyniku}\label{akapit-wyniku}

Podaj N kompletnych par, miarę, CI i jego metodę, odpowiedni test i p,
p\_perm i B. Dla regresji dodaj nachylenie z jednostką i CI, R² oraz
ograniczenie diagnostyczne. Następnie odpowiedz na pytanie merytoryczne
bez zamiany związku w przyczynę.

Kontekst: Hargittai (2002), Internet Skills Scale, DigComp 2.2 i ACRL.
Odsyłacze: \href{../wspolne/literatura.md}{bibliografia kursu}. C08
rozwija odczyt wyników, C09 planuje argument, C10 sprawdza raport i jego
odtwarzalność.

Generatywna SI nie zmienia reguł odczytu: sprawdź nazwę konstruktu,
liczby, N, metodę CI i zgodność czasownika z projektem. Zdanie „p=0,08
oznacza 8\% szans braku efektu'' jest błędne niezależnie od jego autora.
Nie dopisuj brakujących wyników; interpretację piszesz samodzielnie.

\section{Mapa rozwinięć --- bez dodatkowego
zadania}\label{mapa-rozwiniux119ux107-bez-dodatkowego-zadania}

{\def\LTcaptype{none} % do not increment counter
\begin{longtable}[]{@{}
  >{\raggedright\arraybackslash}p{(\linewidth - 4\tabcolsep) * \real{0.3333}}
  >{\raggedright\arraybackslash}p{(\linewidth - 4\tabcolsep) * \real{0.3333}}
  >{\raggedright\arraybackslash}p{(\linewidth - 4\tabcolsep) * \real{0.3333}}@{}}
\toprule\noalign{}
\begin{minipage}[b]{\linewidth}\raggedright
Obraz lub model
\end{minipage} & \begin{minipage}[b]{\linewidth}\raggedright
Co odczytujemy?
\end{minipage} & \begin{minipage}[b]{\linewidth}\raggedright
Główna pułapka
\end{minipage} \\
\midrule\noalign{}
\endhead
\bottomrule\noalign{}
\endlastfoot
\(a+b\exp(-cx)\) & Zmienne tempo spadku średniej czasu & Asymptota nie
jest zmierzonym minimum \\
Wiolina & Gęstość i kształt rozkładu & Szerokość nie musi oznaczać N \\
Pareto & Uporządkowane częstości i kumulację & Najczęstsza bariera nie
musi być najpoważniejsza \\
Koło & Rozłączny podział jednej całości & Wielokrotny wybór osób nie
sumuje się do 100\% \\
Mapa cieplna & Wartość według wspólnej skali koloru & Ten sam procent
może mieć inne N \\
Wydajność & Opóźnienie i przepustowość przy obciążeniu & Czas serwera
nie jest czasem pracy osoby \\
Reszty cząstkowe & \(e_i+b_1x_i\) względem x & Składnik modelu nie jest
całkowitym czasem ani dowodem przyczynowym \\
\end{longtable}
}

Pełny wykład zawiera gotowe wykresy, liczby i komentarze do każdego
przypadku. Galeria służy rozpoznawaniu sposobów prezentacji, nie oznacza
obowiązku stosowania wszystkich metod w projekcie.

\end{document}
